% !TEX root = main.tex
% Full-prefix positional routing (Appendix D). Input after the aligned proof.
% Complete statement and proof of Theorem 4(ii).
\providecommand{\LRnorm}[1]{\left\lVert #1\right\rVert}
\subsection{Full-prefix positional routing}
\label{app:weak-seeded-full-prefix-routing}

All causal sites compete and the current-letter relative bias starts at zero.
The proof has four steps: the exact routing gradient favors the useful site;
a weak feature reaches $x_{\mathrm{boot}}>0$; routing reaches $2/3$ and remains
above $8/13$; a fixed readout comparator completes the loss bound.
The attention floor $A_{\min}$ supports feature growth before routing is learned.

\paragraph{Setting.}\label{app:lr-setting}
Fix $n\ge2$, $p=1/4$, $0<\eta<1/2$ and $0<\varepsilon\le1$.
Use the explicitly tied family \eqref{eq:lr-gates}--\eqref{eq:lr-feature-param}
with $c_0=r_{u,0}=0$, $w_{\mathrm{letter},0}<0$, $v_0,w_{\mathrm{out},0}>0$, and
$w_{\mathrm{state},0}>|b_0|+|w_{\mathrm{letter},0}v_0|$. The fixed query and centered value direction
are as in \cref{app:weak-routing-family}. QK consists only of trainable
positional projections of orthogonal position embeddings; every visible
causal site participates in softmax. All relative biases start at zero.

Assume a complexity bound $\mathsf P\ge2$ bounds $n$, the widths,
$\varepsilon^{-1}$, all initial factor magnitudes, and the Frobenius norms
of the projected QK matrices. Also require
 $\min\{|w_{\mathrm{letter},0}|,v_0,w_{\mathrm{out},0},w_{\mathrm{state},0}-|b_0|-|w_{\mathrm{letter},0}v_0|\}\ge \mathsf P^{-1}.$ 
The initial base scores must have moderate spread: for some $S_0\ge0$,
 $|B_{i_tj}(0)|\le S_0,\qquad e^{2S_0}\le \mathsf P.$ 
Base scores may be bounded perturbations of equal scores.

Our target is \cref{thm:aligned-transformer}(ii):
\begin{equation}\label{eq:lr-finite-theorem}
 \Lstep(T)\le\eta,\qquad T\le\operatorname{poly}(\mathsf P,1/\eta).
\end{equation}

\subsubsection{Exact model and routing gradient}
\label{app:weak-routing-family}

For the fixed query, let $s=(-1)^{s_{t-1}}$. Attention depends only on positions; the pooled centered value is
\begin{equation}\label{eq:lr-centered-value}
 Z_t=\sum_{j\in\mathcal J}A_{tj}(b_j-p),\qquad
 A_{tj}=\frac{\exp(S_{i_tj})}{\sum_{k\le i_t}\exp(S_{i_tk})},
 \quad \mathcal J=\{j_1,\ldots,j_n\}.
\end{equation}
Non-letter sites retain attention mass but have zero value.
A fixed linear direction on orthogonal token/position registers realizes
$b_j-p$ on prompt letters and zero elsewhere; only its amplitude $v$ trains.
For the query token, its match value $1-p$ is canceled by its position
embedding; count/separator tokens and prompt positional corrections are zero.

Let $c,r_u$ be the prior/feature readouts and
$w_{\mathrm{state}},w_{\mathrm{letter}},w_{\mathrm{out}}$ the state input,
letter input and paired output coefficients. The model is
\begin{align}
 r_\pm&=\operatorname{ReLU}(\pm w_{\mathrm{state}}s+w_{\mathrm{letter}}vZ_t+b),
 \label{eq:lr-gates}\\
 f_t&=cs+r_uw_{\mathrm{out}}(r_+-r_-),\label{eq:lr-factor-logit}\\
 i_t-j_t&=n,\qquad t=2,\ldots,n.\label{eq:lr-constant-lag}
\end{align}
The logit is positive for state sign $+1$; the previous state is directly
available in the residual stream. Here $b$ is the MLP bias, distinct from
relative attention biases $b_\ell^{\mathrm{rel}}$.
For $|Z_t|\le1$, the gate condition gives the exact reduction
\begin{align}
 w_{\mathrm{state}}&>|b|+|w_{\mathrm{letter}}v|,\label{eq:lr-gate-condition}\\
 r_+-r_-&=s(w_{\mathrm{state}}+b+w_{\mathrm{letter}}vZ_t),\quad
 f_t=s(\mu-xZ_t),\label{eq:lr-feature-param}\\
 \mu&=c+r_uw_{\mathrm{out}}(w_{\mathrm{state}}+b),\qquad
 x=-r_uw_{\mathrm{out}}w_{\mathrm{letter}}v.\notag
\end{align}

The gate condition activates only $r_+$ for $s=1$ and only $r_-$ for $s=-1$,
which proves the identity. The Euclidean submanifold ties opposite
state/output weights and equal letter weights/biases: duplicated scalars
have metric weight two and gradient factor $1/2$. Absorb the squared norm
of any nonunit fixed matrix direction into its coordinate mobility.

\paragraph{Population loss and routing signs.}
\label{app:weak-routing-signs}

Let $\operatorname{sigmoid}(x)=(1+e^{-x})^{-1}$, $y_t'=1-b_t$ and
\begin{equation}\label{eq:lr-logistic-risk}
 \Lstep=\frac1{n-1}\sum_{t=2}^n\mathbb E
 \left[\log(1+e^{\mu-xZ_t})-y_t'(\mu-xZ_t)\right].
\end{equation}
Multiplying logit and label by $s$ cancels the previous-state sign, giving exactly the binary cross-entropy.
Let $A_t=A_{t,j_t}$ and $\bar A=\frac1{n-1}\sum_{t=2}^n A_t$.

\begin{lemma}[Routing gradient]
\label{lem:lr-signs}\label{lem:lr-postrouting-gradient}
For position-only attention and finite logits,
\begin{align}
 \left.\pd{\Lstep}{x}\right|_{x=0}&=-p(1-p)\bar A<0
 \quad\text{for every }\mu;\label{eq:lr-initial-gate-gradient}\\
 \pd{L_t}{A_{t,j_t}}&<0,\quad \pd{L_t}{A_{tj}}\ge0\ (j\ne j_t),
 \quad \pd{L_t}{A_{tj}}=0\ (j\notin\mathcal J),\label{eq:lr-feature-source-signs}\\
 \pd{L_t}{A_{t,j_t}}&\le-2p(1-p)x\operatorname{sigmoid}(-B),\quad
 \pd{L_t}{A_{tj}}\le p(1-p)x^2A_{tj}/4\ (j\ne j_t).
 \label{eq:lr-source-bounds}
\end{align}
Here $L_t$ is the unaveraged single-position risk; the last two lines
assume $x>0$, and the last also assumes $|\mu|+x\le B$.
For $x>0$ and $A_t\ge1/2$, every competitor score satisfies
\begin{equation}\label{eq:lr-score-comparison}
 \pd{\Lstep}{S_{i_t,j_t}}<0,\qquad
 \pd{\Lstep}{S_{i_tj}}\ge\pd{\Lstep}{S_{i_t,j_t}}.
\end{equation}
Common bias mobility and shared offset $i_t-j_t=n$ therefore give
\begin{equation}\label{eq:lr-bias-gap-mono}
 \frac{d}{dt}(b_n^{\mathrm{rel}}-b_\ell^{\mathrm{rel}})\ge0\quad(\ell\ne n).
\end{equation}
\end{lemma}
\begin{proof}
Independence gives $\mathbb E[(b_t-p)Z_t]=p(1-p)A_t$;
at $x=0$, $\mathbb E Z_t=0$ and
$\pd{\Lstep}{x}=-p(1-p)\bar A$.
For the current letter, condition on all other bits and put
$U=\mu-x\sum_{j\in\mathcal J\setminus\{j_t\}}A_{tj}(b_j-p)$. Then
\[
 \pd{L_t}{A_{t,j_t}}=-xp(1-p)\mathbb E
 \left[1+\operatorname{sigmoid}(U-xA_t(1-p))-\operatorname{sigmoid}(U+xA_tp)\right]<0.
\]
If $|\mu|+x\le B$, the bracket is at least $2\operatorname{sigmoid}(-B)$.
For an incorrect letter $j$, put
$U=\mu-x\sum_{k\in\mathcal J\setminus\{j\}}A_{tk}(b_k-p)$.
The target is independent of $b_j$, so
\[
 \pd{L_t}{A_{tj}}=-xp(1-p)\mathbb E
 \left[\operatorname{sigmoid}(U-xA_{tj}(1-p))-\operatorname{sigmoid}(U+xA_{tj}p)\right].
\]
The sigmoid is increasing and $1/4$-Lipschitz, giving the incorrect-letter
bounds; non-letter value derivatives vanish.

The derivative with respect to a pre-softmax score is
\begin{equation}
 \pd{\Lstep}{S_{i_tj}}=\frac{A_{tj}}{n-1}
 \left(\pd{L_t}{A_{tj}}-\sum_{k\le i_t}A_{tk}\pd{L_t}{A_{tk}}\right).
\end{equation}

Write $D_j=\pd{L_t}{A_{tj}}$ and $\overline D=\sum_k A_{tk}D_k$.
The value signs give $\overline D\ge A_tD_{j_t}$ and $D_{j_t}<\overline D$,
so the correct score derivative is negative.
Since $A_{tj}\le1-A_t\le A_t$, the softmax derivative identity gives
\begin{align*}
 (n-1)\left(\pd{\Lstep}{S_{i_tj}}-\pd{\Lstep}{S_{i_t,j_t}}\right)
 &=A_{tj}D_j-A_tD_{j_t}+(A_t-A_{tj})\overline D\\
 &\ge A_{tj}D_j+A_t(1-A_t+A_{tj})(-D_{j_t})\ge0.
\end{align*}
Offset $n$ occurs in every row and each competitor offset at most once
per row. Summing these comparisons (and the negative correct derivative
where a competitor is absent) gives
$\pd{\Lstep}{b_n^{\mathrm{rel}}}\le\pd{\Lstep}{b_\ell^{\mathrm{rel}}}$,
hence \eqref{eq:lr-bias-gap-mono} under common mobility $m_R>0$.
\end{proof}

\subsubsection{Weak-feature growth}
\label{app:weak-routing-bootstrap}

The fixed positive mobilities are $1$ for $(c,r_u)$, $m_F$ for tied feature
factors, $m_Q$ for positional QK and $m_R$ for all relative biases.

In this subsection, abbreviate $d=w_{\mathrm{state}}+b>0$; $w_{\mathrm{letter}}v<0$.

While the gate regime holds, the tied Euclidean gradient equations include
\begin{align}
 \dot c&=-\pd{\Lstep}{\mu},\qquad
 \dot w_{\mathrm{state}}=\dot b=-\frac{m_Fr_uw_{\mathrm{out}}}{2}\pd{\Lstep}{\mu},
 \label{eq:lr-factor-u}\\
 \dot r_u&=-w_{\mathrm{out}}\left(d\pd{\Lstep}{\mu}-w_{\mathrm{letter}}v\pd{\Lstep}{x}\right),\qquad
 \dot w_{\mathrm{out}}=\frac{m_Fr_u}{2w_{\mathrm{out}}}\dot r_u,\\
 \dot w_{\mathrm{letter}}&=\frac{m_Fr_uw_{\mathrm{out}}v}{2}\pd{\Lstep}{x},\qquad
 \dot v=m_Fr_uw_{\mathrm{out}}w_{\mathrm{letter}}\pd{\Lstep}{x}.
 \label{eq:lr-factor-bv}
\end{align}
In particular,
\begin{equation}\label{eq:lr-factor-invariants}
 r_u^2-\frac{2}{m_F}w_{\mathrm{out}}^2=-\frac{2}{m_F}w_{\mathrm{out},0}^2,
 \qquad
 v^2-2w_{\mathrm{letter}}^2=v_0^2-2w_{\mathrm{letter},0}^2,
 \qquad w_{\mathrm{state}}-b=w_{\mathrm{state},0}-b_0.
\end{equation}
The first invariant gives $w_{\mathrm{out}}^2\ge w_{\mathrm{out},0}^2$.
At $c=c_\star=\log((1-p)/p)$, the gate condition and $r_uw_{\mathrm{out}}\ge0$
make $\dot c\le0$; at $r_u=0$, the velocity points into $r_u>0$.
On $w_{\mathrm{letter}}=0$ or $v=0$,
$\pd{\Lstep}{x}=-p(1-p)\bar A<0$, so the vector field preserves
$w_{\mathrm{letter}}<0,v>0$ (uniqueness excludes simultaneous zero).
Thus $c\le c_\star$, $r_u\ge0$ and $x>0$ persist at positive times
while the gate regime holds.

\begin{lemma}[Feature bootstrap]
\label{lem:lr-feature-bootstrap}
Suppose on an interval prior to reaching the target magnitude that
\begin{equation}\label{eq:lr-amplification-region}
 \bar A\ge A_{\min}>0,\quad
 d\le M_d,\quad
 0<W\le|{w_{\mathrm{letter}} v}|\le W_{\max},\quad
 w_{\mathrm{out}}\ge w_{\mathrm{out},\min}>0,
\end{equation}
and the gate condition \eqref{eq:lr-gate-condition} holds.  Put
\begin{equation}\label{eq:lr-gamma-bootstrap}
 H_F=M_d+W_{\max},\quad
 x_{\mathrm{boot}}=p(1-p) A_{\min}\left(\frac W{H_F}\right)^2,
 \quad
 T_{\mathrm{boot}}=\frac4{3w_{\mathrm{out},\min}^2H_F^2}.
\end{equation}
Starting with $c=r_u=0$, the quantity $x$ reaches $x\ge x_{\mathrm{boot}}$
by $T_{\mathrm{boot}}$, unless one of the region conditions fails first.
\end{lemma}
\begin{proof}
At $c\le c_\star$, since
$0<d+{w_{\mathrm{letter}} v} Z_t\le H_F$, logistic curvature implies
$\pd{\Lstep}{\mu}\le{r_uw_{\mathrm{out}}} H_F/4$.  At $x=0$, $\pd{\Lstep}{x}=-p(1-p)\bar A$ and the
same curvature bound gives $-\pd{\Lstep}{x}\ge p(1-p) A_{\min}-x/4$.
For $x\le x_{\mathrm{boot}}$,
\begin{align*}
 d\pd{\Lstep}{\mu}-w_{\mathrm{letter}} v\pd{\Lstep}{x}
 &\le\frac{M_dH_F}{4W}x-Wp(1-p)A_{\min}+\frac{W_{\max}}4x\\
 &\le-\frac34Wp(1-p)A_{\min}.
\end{align*}
The last step uses
$M_d H_F+WW_{\max}\le H_F^2$.
From \eqref{eq:lr-factor-u} and
$w_{\mathrm{out}}^2\ge w_{\mathrm{out},\min}^2$,
$\ddt{(r_uw_{\mathrm{out}})}\ge(3/4)w_{\mathrm{out},\min}^2Wp(1-p) A_{\min}$.
In this region $\pd{\Lstep}{x}<0$ and
$\ddt{(w_{\mathrm{letter}} v)}=m_F{r_uw_{\mathrm{out}}}(v^2+w_{\mathrm{letter}}^2/2)\pd{\Lstep}{x}\le0$.
Therefore
$\dot x\ge(3/4)w_{\mathrm{out},\min}^2W^2p(1-p)A_{\min}$.
Dividing $x_{\mathrm{boot}}$ by this rate gives exactly
$4/(3w_{\mathrm{out},\min}^2H_F^2)=T_{\mathrm{boot}}$.
\end{proof}

The exact metric weights are $(2,2,2,2,1)$ for
$F=(w_{\mathrm{state}},b,w_{\mathrm{letter}},w_{\mathrm{out}},v)$ and
$\varepsilon^{-2}$ for projected positional rows $(Q,K)$.
Orthogonal embeddings $\varepsilon e_i$ give
$\dot Q=-m_Q\varepsilon^2\nabla_Q\Lstep$,
$\dot K=-m_Q\varepsilon^2\nabla_K\Lstep$, hence the loss-energy identity
\begin{equation}\label{eq:lr-mobility-budget}
 -\ddt{\Lstep}=
 \LRnorm{\nabla_{c,r_u}\Lstep}^2
 +m_F\LRnorm{\nabla_F\Lstep}^2
 +m_Q\LRnorm{\nabla_{Q,K}\Lstep}^2
 +m_R\LRnorm{\nabla_{b^{\mathrm{rel}}}\Lstep}^2.
\end{equation}
Here group norms and gradients use those induced metrics. For each group
with mobility $m_G$, Cauchy--Schwarz implies
\begin{equation}\label{eq:lr-movement-budget}
 \LRnorm{\theta_G(T)-\theta_G(0)}
 \le\sqrt{m_G T\,[\Lstep(0)-\Lstep(T)]}
 \le\sqrt{m_G T\log2}.
\end{equation}


\subsubsection{Routing amplification, retention and loss completion}
\label{app:weak-routing-convergence}

Use positional scores
\begin{equation}\label{eq:lr-relative-scores}
 S_{i_tj}=B_{i_tj}+b_{i_t-j}^{\mathrm{rel}},
 \qquad
 B_{i_tj}=\frac{q_{i_t}^\top k_j}{\sqrt{d_h}},
 \qquad b_\ell^{\mathrm{rel}}(0)=0.
\end{equation}

\begin{lemma}[Routing acquisition and retention]
\label{lem:lr-bias-learning}\label{lem:lr-post-routing}
With $|B_{i_tj}|\le S_A$, zero initial relative biases and
$0<x\le A_{\min}/4$, the condition $A_t\ge A_{\min}>0$ implies
\begin{align}
 \pd{\Lstep}{S_{i_t,j_t}}&<0,\quad \pd{\Lstep}{S_{i_tj}}\ge0\ (j\ne j_t),
 \label{eq:lr-smallscore-sign}\\
 \dot b_n^{\mathrm{rel}}&>0,\quad \dot b_\ell^{\mathrm{rel}}\le0\ (\ell\ne n),
 \label{eq:lr-beta-routes}\\
 A_t&\ge A_{\min}:=e^{-2S_A}/m,\quad
 m:=\max_t\#\{j:j\le i_t\}\le3n.\label{eq:lr-floor}
\end{align}
After $x=x_{\mathrm{boot}}<A_{\min}/4$, a polynomially bounded fixed $m_R$
brings every $A_t$ to $2/3$ before $x$ leaves
$[x_{\mathrm{boot}}/2,A_{\min}/4]$, provided the feature/base-score
neighbourhoods hold. Once $A_t(t_0)\ge2/3$, if $x>0$, the feature family
persists and
$|[B_{i_t,j_t}-B_{i_tj}](t)-[B_{i_t,j_t}-B_{i_tj}](t_0)|\le\log(5/4)$,
then
\begin{equation}\label{eq:lr-post-floor}
 A_t(t)\ge8/13>3/5.
\end{equation}
\end{lemma}
\begin{proof}
\emph{Acquisition.} By Lemma~\ref{lem:lr-signs}, the correct value derivative is negative;
incorrect-letter derivatives are nonnegative.  The correct derivative has the logit-independent bound
 $-\pd{L_t}{A_{t,j_t}}\ge p(1-p) x(1-xA_t/4)\ge p(1-p) x(1-x/4),$ 
because $\operatorname{sigmoid}$ is $1/4$-Lipschitz. Hence
 $\sum_{k\le i_t}A_{tk}\pd{L_t}{A_{tk}}\le-p(1-p) x A_t(1-x/4)+p(1-p) x^2/4<0$ 
when $x\le A_{\min}/4$ and $A_t\ge A_{\min}$.  Hence the softmax derivatives have the stated signs.
The bias sign conclusion follows from the fixed offset
\eqref{eq:lr-constant-lag}.
With $b_n^{\mathrm{rel}}\ge0$, all other biases $\le0$,
and $|B|\le S_A$, softmax normalization yields
$A_t\ge e^{-2S_A}/m$.

If routing is incomplete, some row has $A_t\le2/3$. In that row,
\[
 -\pd{\Lstep}{S_{i_t,j_t}}
 \ge\frac{A_t(1-A_t)}{n-1}(-\pd{L_t}{A_{t,j_t}})
 \ge\frac{p(1-p)xA_{\min}}{4(n-1)}.
\]
For $x\ge x_{\mathrm{boot}}/2$, summing the nonnegative correct-score
credits yields $\dot b_n^{\mathrm{rel}}\ge
m_Rp(1-p)x_{\mathrm{boot}}A_{\min}/[8(n-1)]$.
Every row has $A_t\ge2/3$ once
$b_n^{\mathrm{rel}}\ge2S_A+\log(2(m-1))$, so
\[
 T_{\mathrm{route}}\le
 \frac{8(n-1)[2S_A+\log(2(m-1))]}{m_Rp(1-p)x_{\mathrm{boot}}A_{\min}}.
\]
The constructive proof below chooses $m_R$ from a uniform bound
$|\dot x|\le M_x$ to keep this time within the small-feature region.
For $p=1/4$, $x_{\mathrm{boot}}\le3A_{\min}/16<A_{\min}/4$.

\emph{Retention.} While every $A_t\ge1/2$, \cref{lem:lr-signs} makes
all relative-bias gaps nondecreasing. Base-score drift multiplies each
competitor's odds by at most $5/4$. Their sum was at most $1/2$ at $t_0$,
hence remains at most $5/8$: $A_t\ge1/(1+5/8)=8/13$.
This excludes a first exit through $A_t=1/2$.
\end{proof}

\begin{lemma}[Loss completion]
\label{lem:lr-readout-convergence}
If $A_t\ge3/5$ at every supervised position, every prompt realization obeys
\begin{equation}\label{eq:lr-bit-separation}
 (1-2b_t)(1/2-p-Z_t)\ge1/10.
\end{equation}
Suppose the feature neighbourhood admits a fixed comparator
$\widehat r=(\widehat c,\widehat r_u)$ throughout $[t_0,t_0+T]$, with
\begin{equation}\label{eq:lr-comparator}
 \Lstep(\widehat r;F(t))\le\eta/2.
\end{equation}
With unit readout mobility and nonnegative mobility for every other group,
\begin{equation}\label{eq:lr-descent}
 \Lstep(t_0+T)\le\eta/2+\LRnorm{r(t_0)-\widehat r}^2/(2T).
\end{equation}
Thus $T=\LRnorm{r(t_0)-\widehat r}^2/\eta$ suffices for $\Lstep\le\eta$.
\end{lemma}
\begin{proof}
For $b_t=0$, $Z_t\le1-p-A_t$; for $b_t=1$,
$Z_t\ge A_t-p$.  Both cases imply
\eqref{eq:lr-bit-separation}.
At the seed, put $d_0=w_{\mathrm{state},0}+b_0$ and fix
\[
 \widehat r_u=\frac{20\log(2/\eta)}{w_{\mathrm{out},0}|w_{\mathrm{letter},0}v_0|},\qquad
 \widehat c=20\log(2/\eta)(1/2-p)-\widehat r_uw_{\mathrm{out},0}d_0.
\]
At the reference feature, its effective logit is
$20\log(2/\eta)(1/2-p-Z_t)$, which has signed margin at least $2\log(2/\eta)$.
Require the feature neighbourhood to satisfy
\begin{equation}\label{eq:lr-comparator-stability}
 |w_{\mathrm{out}}(w_{\mathrm{state}}+b)-w_{\mathrm{out},0}d_0|+|w_{\mathrm{out}} {w_{\mathrm{letter}} v}-w_{\mathrm{out},0}{w_{\mathrm{letter},0}v_0}|
 \le\frac{w_{\mathrm{out},0}|{w_{\mathrm{letter},0}v_0}|}{20}.
\end{equation}
Since $|Z_t|\le1$, the perturbation costs at most $\log(2/\eta)$ of margin,
leaving margin at least $\log(2/\eta)$ and loss at most $\eta/2$.
This comparator depends only on the seed, not the routing hitting time.
Convexity in $r=(c,r_u)$ \citep[Sections~3.1.3, 3.1.5]{boyd2004convex} gives
$\frac12\frac{d}{dt}\LRnorm{r-\widehat r}^2
\le\Lstep(\widehat r;F(t))-\Lstep(r;F(t))$.
Integrate and use nonincreasing joint-flow loss to obtain \eqref{eq:lr-descent}.
\end{proof}

\begin{proof}[Proof of \cref{thm:aligned-transformer}(ii)]
Choose initial coordinate balls with inverse-polynomial radii $r_F,r_Q$.
The feature ball preserves \eqref{eq:lr-gate-condition},
\eqref{eq:lr-comparator-stability}, $w_{\mathrm{letter}}<0,v>0$ and the
following bounds; the QK ball preserves $|B|\le S_A=S_0+1$ and each
base-score contrast within $\frac12\log(5/4)$ of its initial value.
\begin{center}
\small
\begin{tabular}{@{}ll@{}}
\toprule
Feature quantity & Uniform bound \\
\midrule
$|w_{\mathrm{letter}}v|$ & $W\le|w_{\mathrm{letter}}v|\le W_{\max}$ \\
$w_{\mathrm{state}}+b$ & $w_{\mathrm{state}}+b\le M_d$, $H_F=M_d+W_{\max}$ \\
$w_{\mathrm{out}}$ & $w_{\mathrm{out},\min}\le w_{\mathrm{out}}\le M_{\mathrm{out}}$ \\
\bottomrule
\end{tabular}
\end{center}
Seed margins and polynomial Lipschitz bounds give these radii and polynomial
upper bounds and reciprocal positive lower bounds. Any two QK points
in the ball have contrast differences at most $\log(5/4)$.

All choices below depend only on these initial bounds, in the displayed
order. Set $A_{\min}=e^{-2S_A}/m$ and use
$x_{\mathrm{boot}},T_{\mathrm{boot}}$ from \eqref{eq:lr-gamma-bootstrap}.
Before $T_{\mathrm{pre}}=T_{\mathrm{boot}}+1$, the readout equations give
$|\dot c|\le1$, $|\dot r_u|\le M_{\mathrm{out}}H_F$, hence
$\|r(t)\|\le T_{\mathrm{pre}}\sqrt{1+M_{\mathrm{out}}^2H_F^2}$.
These estimates use $|\operatorname{sigmoid}(\mu-xZ_t)-y_t'|\le1$ and
$|Z_t|\le1$. With $m_F\le1$, the factor equations therefore give a
polynomial $M_x\ge1$ bounding $|\dot x|$, uniformly over attention and
independently of $m_R,m_Q$. Use the comparator of
\cref{lem:lr-readout-convergence} and fix:
\begin{center}
\small
\begin{tabular}{@{}p{.20\linewidth}p{.74\linewidth}@{}}
\toprule
Choice & Definition and guarantee \\
\midrule
Routing window & $T_{\max}=\min\{1,x_{\mathrm{boot}}/(4M_x),
(A_{\min}/4-x_{\mathrm{boot}})/(2M_x)\}$;
$K_R=2S_A+\log(2(m-1))$, $c_0=1/4$. \\
Relative biases & $m_R=\max\{1,2(n-1)K_R/\bigl(c_0p(1-p)x_{\mathrm{boot}}A_{\min}T_{\max}\bigr)\}$;
routing takes at most $T_{\max}$. \\
Completion horizon & $R_{\mathrm{cmp}}=T_{\mathrm{pre}}\sqrt{1+M_{\mathrm{out}}^2H_F^2}+\|\widehat r\|$;
$T_{\mathrm{post}}=R_{\mathrm{cmp}}^2/\eta$,
$T^\star=T_{\mathrm{pre}}+T_{\mathrm{post}}$. \\
Feature factors & $m_F=\min\{1,r_F^2/(4T^\star\log2)\}$;
feature movement is at most $r_F/2$. \\
Positional QK & $m_Q=\min\{1,r_Q^2/(4T^\star\log2)\}$;
QK movement is at most $r_Q/2$. \\
Readouts & Unit mobility throughout; convex loss completion remains active. \\
\bottomrule
\end{tabular}
\end{center}
All bounds and relevant inverses are polynomial in $\mathsf P,1/\eta$;
$x_{\mathrm{boot}}\le3A_{\min}/16$ makes $T_{\max}>0$. The uniform $M_x$
was chosen for every $m_F\le1$, before fixing $m_F$.

Metric weights are at least one, so \eqref{eq:lr-movement-budget}
limits coordinate movement to half each ball radius up to $T^\star$.
A first exit is impossible. Smoothness in the gate region and the energy
budget also exclude finite-time escape of any group. The inward signs
preserve $c\le c_\star$, $r_u\ge0$ and $x>0$ at positive times.

Starting from zero biases, \cref{lem:lr-bias-learning} preserves
$A_t\ge A_{\min}$ until \cref{lem:lr-feature-bootstrap} gives
$x=x_{\mathrm{boot}}$ by $T_{\mathrm{boot}}$.
The bound $|\dot x|\le M_x$ then keeps $x\in[x_{\mathrm{boot}}/2,A_{\min}/4]$
during acquisition, which takes at most $T_{\max}$ (zero if already routed).
Thus some $t_0\le T_{\mathrm{pre}}$ has every $A_t\ge2/3$;
retention keeps $A_t\ge8/13>3/5$ afterwards.
The feature ball preserves the comparator and
$\|r(t_0)-\widehat r\|\le R_{\mathrm{cmp}}$, so
\cref{lem:lr-readout-convergence} gives
$\Lstep(t_0+T_{\mathrm{post}})\le\eta$ by $T^\star$.
All mobilities depend only on initial bounds.
\end{proof}

\begin{corollary}[Conditional autoregressive rollout]
\label{cor:conditional-rollout}
If first-state generation, formatting and final copying are exact,
sampled decoding with teacher-forced transition loss $\Lstep\le\eta$
has trace/answer error probability at most $\min\{1,(n-1)\eta\}$.
\end{corollary}
\begin{proof}
On a correct prefix, the error probability is $1-q_t\le-\log q_t$, where
$q_t$ is the correct-state probability. Sum over first-error positions and
average over prompts to obtain $(n-1)\Lstep$; the exact boundary operations
add no errors.
\end{proof}
